
\newpage
\section*{Rychlé opakování před zkouškou}
\addcontentsline{toc}{section}{Rychlé opakování před zkouškou}
\begin{itemize}
  \item Hrozba není riziko. Hrozba je zdroj škody, riziko je míra ohrožení konkrétního objektu.
  \item Kyberbezpečnost stírá hranici vnitřní a vnější bezpečnosti.
  \item Strategie není zákon. Strategie určuje cíle, zákon ukládá povinnosti, akční plán rozepisuje úkoly.
  \item Ne každý kyberútok je kyberválka. Kyberválku říkej jen u operací v ozbrojeném konfliktu nebo s účinky srovnatelnými se závažným vojenským útokem.
  \item U kyberválky měj připravené příklady: APT31 = špionáž a politická atribuce, Volt Typhoon = pre-positioning, Salt Typhoon = telekomunikace, APT28/TP-Link = cizí infrastruktura, NoName057(16) = hacktivismus/netwar, Viasat = souběh s invazí.
  \item U kritické infrastruktury nejde jen o data, ale o kontinuitu základní služby a fyzické dopady.
  \item CER není český zákon ani "nařízení CER"; směrnice CER je v ČR transponována hlavně zákonem č. 266/2025 Sb. a praktické určení subjektů závisí na účinném nařízení vlády o základních službách a kritériích významnosti.
  \item Starý systém KI byl objektový, nový systém je orientovaný na základní službu a poskytovatele; subjekt KI vzniká až doručením rozhodnutí Ministerstva vnitra.
  \item NIS2 se v ČR učí přes zákon o kybernetické bezpečnosti a prováděcí vyhlášky.
  \item U kyberkriminality převáděj technický popis na skutkové podstaty.
  \item U elektronických důkazů hlídej volatilitu, integritu, hash a chain of custody.
  \item U GDPR nejdřív účel, potom právní titul, minimalizace, riziko a dokumentace.
  \item Datová schránka není elektronický podpis, ale zákon jí může dát účinky podepsaného úkonu.
\end{itemize}

\section*{Nově doplněné příklady podle otázek}
\addcontentsline{toc}{section}{Nově doplněné příklady podle otázek}
\begin{itemize}
  \item \textbf{15}: APT31 proti infrastruktuře Ministerstva zahraničních věcí ČR jako průřezový příklad referenčního objektu, hrozby, rizika, vnitřní/vnější bezpečnosti, sektorů a aktérů bezpečnostní politiky.
  \item \textbf{16}: APT31 jako řetězec od technického zjištění po veřejnou atribuci vlády; DDoS vlna v Q1 2026 jako příklad rozdílu mezi provozním incidentem a strategickou reakcí státu.
  \item \textbf{17}: APT31, Volt Typhoon, Salt Typhoon, APT28/TP-Link, NoName057(16), Viasat, Ukrajina, AI, dodavatelské řetězce a soukromý sektor jako konkrétní příklady kyberkonfliktu pod prahem i během ozbrojeného konfliktu.
  \item \textbf{18}: přechod starý/nový systém KI; Synnovis/NHS, Volt Typhoon, Viasat, ukrajinská energetika 2015/2016, Colonial Pipeline a NoName057(16) jako sada příkladů pro základní službu, IT/OT, dodavatele, kontinuitu a netwar.
  \item \textbf{19}: modelové situace nemocnice, banky, cloudového dodavatele, výrobce routeru podle CRA a použití Cyber Solidarity Act.
  \item \textbf{20}: Operation Endgame, First VPN, Cryptomixer/AudiA6 a modelové kvalifikace pro infostealery, deepfake podvod, initial access brokera, DDoS-for-hire a kryptoměnové praní.
  \item \textbf{21}: praktické použití § 78, § 79, § 82, § 7b, § 113a, § 88a, § 88, znaleckého zkoumání, přeshraničního scénáře a forenzní osy u ransomwaru.
  \item \textbf{22}: GDPR přes účel, právní titul, zásady, zvláštní kategorie, oprávněný zájem, DPIA, DPO, práva subjektů a ÚOOÚ.
  \item \textbf{23}: elektronický podpis, uznávaný podpis, QES, pečeť, časové razítko, datové schránky a autorizovaná konverze.
\end{itemize}

\section*{Skutečné a modelové příklady}
\addcontentsline{toc}{section}{Skutečné a modelové příklady}
\begin{itemize}
  \item \textbf{Skutečné případy z oficiálních zdrojů}: APT31/MZV ČR, DDoS vlna Q1 2026, Volt Typhoon, Salt Typhoon, APT28/TP-Link, NoName057(16), Viasat, Colonial Pipeline, ukrajinská kybernetická obrana podle CCDCOE, Synnovis/NHS, Operation Endgame, First VPN a Cryptomixer/AudiA6.
  \item \textbf{Modelové příklady pro vysvětlení práva}: vodárenská společnost, nemocnice a ransomware podle ZKB/NIS2, banka a cloud, výrobce routeru podle CRA, e-evidence scénář, použití procesních institutů trestního řádu, phishing/cloud přes hranice, GDPR kamera/sklad, QES, pečeť, časové razítko a datové schránky.
  \item \textbf{Atribuce}: u státních aktérů rozlišuj veřejnou politickou atribuci státu, technické přiřazení bezpečnostní komunitou, pravděpodobné spojení se státem a právní přičitatelnost. Ne každý případ kyberšpionáže je kyberválka a ne každá zdrojová IP adresa identifikuje pachatele.
\end{itemize}

\section*{Dynamické údaje k ověření před státnicí}
\addcontentsline{toc}{section}{Dynamické údaje k ověření před státnicí}
\begin{itemize}
  \item \textbf{E-evidence}: od 18. 8. 2026 se začne používat EPOC/EPOC-PR; před zkouškou ověř, zda už je režim použitelný, jak je nastaven český příjem a validace příkazů a zda se změnila praxe poskytovatelů.
  \item \textbf{2. dodatkový protokol k Budapešťské úmluvě}: ověř počet ratifikací a zda už vstoupil v platnost; k 18. 6. 2026 mu chyběla pátá ratifikace.
  \item \textbf{Úmluva OSN proti kyberkriminalitě}: ověř počet signatářů, ratifikací a vstup v platnost; k 18. 6. 2026 nebyla v platnosti a UNTC uváděl 77 signatářů a 3 strany.
  \item \textbf{CER/ZKI}: k 18. 6. 2026 nebylo v dostupných oficiálních zdrojích potvrzeno účinné nařízení vlády o základních službách a kritériích významnosti ani první individuální rozhodnutí MV o zařazení nových subjektů. Těsně před státnicí ověř e-Sbírku, MV a GŘ HZS ČR: zda už je nařízení účinné, zda byla vydána rozhodnutí, zda se změnil termín 17. 7. 2026 a zda existuje novela nebo nové stanovisko k přechodu.
  \item \textbf{ZKB/NIS2}: ověř aktuální prováděcí předpisy NÚKIB, metodiky, režim strategicky významných služeb a pravidla bezpečnosti dodavatelského řetězce.
  \item \textbf{eIDAS 2.0 a EUDI Wallet}: ověř prováděcí akty, český projekt DIA, povinnosti přijímat peněženku a přesné lhůty pro zpřístupnění.
  \item \textbf{Předávání osobních údajů do USA}: ověř stav EU-US Data Privacy Framework a případná rozhodnutí Soudního dvora EU navazující na věc Latombe.
  \item \textbf{Strategie}: ověř aktuální akční plán Národní strategie kybernetické bezpečnosti 2026-2030 a dokumenty ke Strategii pro posílení odolnosti subjektů kritické infrastruktury 2026-2029.
\end{itemize}

\section*{Ověřené online zdroje}
\addcontentsline{toc}{section}{Ověřené online zdroje}
\begin{itemize}
  \item NÚKIB - strategie a akční plány: \url{https://nukib.gov.cz/cs/infoservis/dokumenty-a-publikace/strategie-akcni-plan/}
  \item NÚKIB - veřejná atribuce kampaně APT31 proti MZV ČR: \url{https://nukib.gov.cz/cs/infoservis/aktuality/2262-ceska-vlada-provedla-verejnou-atribuci-kybernetickych-utoku-ciny-apt31-napojena-na-cinskou-zpravodajskou-sluzbu-dlouhodobe-cilila-na-infrastrukturu-ministerstva-zahranicnich-veci-cr/}
  \item Ministerstvo zahraničních věcí - prohlášení vlády k APT31: \url{https://mzv.gov.cz/jnp/cz/udalosti_a_media/tiskove_zpravy/prohlaseni_vlady_ceske_republiky.html}
  \item NÚKIB - čtvrtletní přehled hrozeb Q1 2026: \url{https://nukib.gov.cz/cs/infoservis/aktuality/2399-vydali-jsme-ctvrtletni-prehled-hrozeb-nukib-q1-2026/}
  \item NÚKIB - Zpráva o stavu kybernetické bezpečnosti ČR za rok 2024: \url{https://nukib.gov.cz/download/publikace/zpravy_o_stavu/Zpr%C3%A1va_o_stavu_kyberneticke_bezpecnosti_CR_za_rok_2024.pdf}
  \item NÚKIB - routery TP-Link kompromitované aktérem APT28: \url{https://nukib.gov.cz/cs/infoservis/aktuality/2394-routery-tp-link-kompromitovany-ruskym-statnim-akterem-apt28/}
  \item NÚKIB - proruské hacktivistické skupiny proti kritické infrastruktuře: \url{https://nukib.gov.cz/cs/infoservis/aktuality/2356-nukib-se-pripojil-k-mezinarodnimu-upozorneni-na-proruske-hacktivisticke-skupiny-utocici-na-subjekty-kriticke-infrastruktury/}
  \item NÚKIB - Salt Typhoon a kompromitace sítí: \url{https://nukib.gov.cz/cs/infoservis/aktuality/2292-nukib-spolu-s-nsa-a-dalsimi-americkymi-urady-upozornuje-na-cinskeho-aktera-salt-typhoon-ktery-kompromituje-site-po-celem-svete/}
  \item Portál NÚKIB k novému zákonu o kybernetické bezpečnosti: \url{https://portal.nukib.gov.cz/}
  \item e-Sbírka - zákon o kybernetické bezpečnosti (č. 264/2025 Sb.): \url{https://e-sbirka.gov.cz/sb/2025/264}
  \item e-Sbírka - zákon o odolnosti subjektů kritické infrastruktury (č. 266/2025 Sb.): \url{https://e-sbirka.gov.cz/sb/2025/266}
  \item HZS ČR - kritická infrastruktura po nové úpravě: \url{https://hzscr.gov.cz/clanek/web-informacni-servis-zpravodajstvi-2026-unor-kriticka-infrastruktura.aspx}
  \item HZS ČR - představení zákona o kritické infrastruktuře: \url{https://hzscr.gov.cz/clanek/predstaveni-zakona-o-kriticke-infrastrukture.aspx}
  \item GŘ HZS ČR - metodika k zařazení poskytovatele základní služby mezi subjekty KI: \url{https://hzscr.gov.cz/soubor/2026-02-25-metodika-procesu-urceni-subjektu-ki-pro-subjekty-pdf.aspx}
  \item Ministerstvo vnitra - poskytování informací podle zákona o kritické infrastruktuře: \url{https://mv.gov.cz/npo/clanek/sdeleni-ministerstva-vnitra-k-poskytovani-informaci-podle-zakona-o-kriticke-infrastrukture.aspx}
  \item Vláda ČR - Bezpečnostní rada státu: \url{https://vlada.gov.cz/cz/pracovni-a-poradni-organy-vlady/brs/brs-uvod-3851/}
  \item EUR-Lex - NIS2: \url{https://eur-lex.europa.eu/eli/dir/2022/2555/oj}
  \item EUR-Lex - CER: \url{https://eur-lex.europa.eu/eli/dir/2022/2557/oj}
  \item EUR-Lex - Cybersecurity Act 2019/881: \url{https://eur-lex.europa.eu/eli/reg/2019/881/oj}
  \item EUR-Lex - Cyber Solidarity Act 2025/38: \url{https://eur-lex.europa.eu/eli/reg/2025/38/oj}
  \item EUR-Lex - útoky na informační systémy, směrnice 2013/40/EU: \url{https://eur-lex.europa.eu/eli/dir/2013/40/oj}
  \item EUR-Lex - bezhotovostní platební podvody, směrnice 2019/713: \url{https://eur-lex.europa.eu/eli/dir/2019/713/oj}
  \item EUR-Lex - evropský vyšetřovací příkaz, směrnice 2014/41/EU: \url{https://eur-lex.europa.eu/eli/dir/2014/41/oj}
  \item EUR-Lex - e-evidence nařízení 2023/1543: \url{https://eur-lex.europa.eu/eli/reg/2023/1543/oj}
  \item EUR-Lex - e-evidence směrnice 2023/1544: \url{https://eur-lex.europa.eu/eli/dir/2023/1544/oj}
  \item EUR-Lex - Cyber Resilience Act 2024/2847: \url{https://eur-lex.europa.eu/eli/reg/2024/2847/oj}
  \item EUR-Lex - eIDAS 2.0, nařízení 2024/1183: \url{https://eur-lex.europa.eu/eli/reg/2024/1183/oj}
  \item EUR-Lex - GDPR, nařízení 2016/679: \url{https://eur-lex.europa.eu/eli/reg/2016/679/oj}
  \item EUR-Lex - policejní směrnice 2016/680: \url{https://eur-lex.europa.eu/eli/dir/2016/680/oj}
  \item NATO CCDCOE - analýza kyberútoků proti Estonsku 2007: \url{https://ccdcoe.org/uploads/2018/10/Ottis2008_AnalysisOf2007FromTheInformationWarfarePerspective.pdf}
  \item CISA - kyberútok proti ukrajinské kritické infrastruktuře 2015: \url{https://www.cisa.gov/news-events/ics-alerts/ir-alert-h-16-056-01}
  \item CISA - Volt Typhoon a kritická infrastruktura: \url{https://www.cisa.gov/news-events/cybersecurity-advisories/aa24-038a}
  \item CISA - Volt Typhoon kompromitace SOHO zařízení: \url{https://www.cisa.gov/news-events/cybersecurity-advisories/aa23-144a}
  \item CISA - Salt Typhoon / čínští státní aktéři proti telekomunikacím: \url{https://www.cisa.gov/news-events/cybersecurity-advisories/aa25-239a}
  \item CISA - Colonial Pipeline: \url{https://www.cisa.gov/news-events/news/attack-colonial-pipeline-what-weve-learned-what-weve-done-over-past-two-years}
  \item U.S. Department of Energy - Colonial Pipeline Cyber Incident: \url{https://www.energy.gov/ceser/colonial-pipeline-cyber-incident}
  \item Rada EU - ruské kybernetické operace proti Ukrajině a KA-SAT: \url{https://www.consilium.europa.eu/en/press/press-releases/2022/05/10/russian-cyber-operations-against-ukraine-declaration-by-the-high-representative-on-behalf-of-the-european-union/}
  \item Viasat - KA-SAT Network cyber attack overview: \url{https://www.viasat.com/perspectives/corporate/2022/ka-sat-network-cyber-attack-overview/}
  \item UK Government - atribuce NotPetya Rusku: \url{https://www.gov.uk/government/news/foreign-office-minister-condemns-russia-for-notpetya-attacks}
  \item White House Archives - NotPetya attribution statement: \url{https://trumpwhitehouse.archives.gov/briefings-statements/statement-press-secretary-25/}
  \item UK NCSC - APT28 a zneužití routerů k DNS hijackingu: \url{https://www.ncsc.gov.uk/news/apt28-exploit-routers-to-enable-dns-hijacking-operations}
  \item NATO CCDCOE - ukrajinská kybernetická obrana a role nestátních aktérů: \url{https://ccdcoe.org/library/publications/policy-brief-ukraines-cyber-defence-evolution-the-role-of-non-state-actors-and-public-private-partnerships/}
  \item ENISA - Space Threat Landscape 2025: \url{https://www.enisa.europa.eu/publications/enisa-space-threat-landscape-2025}
  \item NHS England - Synnovis cyber attack statement: \url{https://www.england.nhs.uk/2024/06/synnovis-cyber-attack-statement-from-nhs-england/}
  \item NHS England - Synnovis cyber incident: \url{https://www.england.nhs.uk/synnovis-cyber-incident/}
  \item Europol - IOCTA 2026: \url{https://www.europol.europa.eu/publications-events/main-reports/iocta-report}
  \item Europol - Operation Endgame 2025: \url{https://www.europol.europa.eu/media-press/newsroom/news/end-of-game-for-cybercrime-infrastructure-1025-servers-taken-down}
  \item Eurojust - Operation Endgame 2025: \url{https://www.eurojust.europa.eu/news/authorities-continue-protect-citizens-cybercriminals-during-major-malware-operation}
  \item Eurojust - First VPN: \url{https://www.eurojust.europa.eu/news/eurojust-coordinated-investigation-shuts-down-criminal-vpn-network}
  \item Europol - First VPN: \url{https://www.europol.europa.eu/media-press/newsroom/news/cybercriminal-vpn-used-ransomware-actors-dismantled-in-global-crackdown}
  \item Eurojust - Cryptomixer: \url{https://www.eurojust.europa.eu/news/cryptocurrency-mixing-service-used-launder-money-taken-down}
  \item Europol - AudiA6 crypto laundering pipeline: \url{https://www.europol.europa.eu/media-press/newsroom/news/ransomware-gangs-cut-eur-336-million-audia6-crypto-laundering-pipeline}
  \item ÚOOÚ - posouzení vlivu na ochranu osobních údajů: \url{https://uoou.gov.cz/profesional/metodiky-a-doporuceni-pro-spravce/posouzeni-vlivu-na-ochranu-osobnich-udaju}
  \item ÚOOÚ - základní příručka k ochraně údajů: \url{https://uoou.gov.cz/verejnost/zakladni-prirucka-k-ochrane-udaju}
  \item ÚOOÚ - kontrola biometrické docházky ve výrobní společnosti: \url{https://uoou.gov.cz/cinnost/ochrana-osobnich-udaju/ukoncene-kontroly/kontroly-za-rok-2025/kontrolni-cinnost-v-oblasti-ochrany-osobnich-udaju-2025/vyrobni-spolecnost-uoou-0433023}
  \item ÚOOÚ - realitní zprostředkovatel a následné využití údajů: \url{https://uoou.gov.cz/cinnost/ochrana-osobnich-udaju/ukoncene-kontroly/kontroly-za-rok-2025/kontrolni-cinnost-v-oblasti-ochrany-osobnich-udaju-2025/realitni-zprostredkovatel-uoou-0351021}
  \item ÚOOÚ - investiční zprostředkovatel a právní titul zpracování: \url{https://uoou.gov.cz/cinnost/ochrana-osobnich-udaju/ukoncene-kontroly/kontroly-za-rok-2025/kontrolni-cinnost-v-oblasti-ochrany-osobnich-udaju-2025/investicni-zprostredkovatel-uoou-0248922}
  \item EDPB - územní působnost GDPR: \url{https://www.edpb.europa.eu/our-work-tools/our-documents/guidelines/guidelines-32018-territorial-scope-gdpr-article-3-version_en}
  \item EDPB - Guidelines 1/2024 k oprávněnému zájmu: \url{https://www.edpb.europa.eu/our-work-tools/documents/public-consultations/2024/guidelines-12024-processing-personal-data-based_en}
  \item EDPB - stanovisko 28/2024 k AI modelům a GDPR: \url{https://www.edpb.europa.eu/our-work-tools/our-documents/opinion-board-art-64/opinion-282024-certain-data-protection-aspects_en}
  \item ÚOOÚ - Avast/Jumpshot pokuta 351 mil. Kč: \url{https://uoou.gov.cz/novinky/vse/uoou-ulozil-pokutu-351-mil-kc-za-poruseni-gdpr}
  \item Městský soud v Praze - Avast/Jumpshot rozsudek 2025: \url{https://msp.gov.cz/documents/d/mestsky-soud-v-praze/tiskova-zprava-k-rozsudku-ze-dne-30-9-2025-sp-zn-5-a-56-2024}
  \item CURIA - Latombe v Commission, T-553/23: \url{https://infocuria.curia.europa.eu/tabs/redirect/juris/liste.jsf?num=T-553/23}
  \item Council of Europe - Budapešťská úmluva a protokoly: \url{https://www.coe.int/en/web/cybercrime/the-budapest-convention}
  \item Council of Europe - 2. dodatkový protokol CETS 224, podpisy a ratifikace: \url{https://www.coe.int/en/web/conventions/full-list?module=signatures-by-treaty&treatynum=224}
  \item UN Treaty Collection - Úmluva OSN proti kyberkriminalitě: \url{https://treaties.un.org/pages/ViewDetails.aspx?chapter=18&clang=_en&mtdsg_no=XVIII-16&src=TREATY}
  \item DIA - evropská peněženka digitální identity: \url{https://www.dia.gov.cz/eudiw/cs}
  \item DIA - seznam kvalifikovaných poskytovatelů služeb vytvářejících důvěru: \url{https://www.dia.gov.cz/cs/legislativa/eidas-sluzby-vytvarejici-duveru-a-elektronicka-identifikace/povinne-zverejnovane-informace/seznam-kvalifikovanych-poskytovatelu-sluzeb-vytvarejicich-duveru-a-poskytovanych-kvalifikovanych-sluzeb-vytvarejicich-duveru}
  \item DIA - seznam kvalifikovaných prostředků pro vytváření elektronických podpisů v ČR: \url{https://www.dia.gov.cz/cs/legislativa/eidas-sluzby-vytvarejici-duveru-a-elektronicka-identifikace/informace-pro-uzivatele/seznam-vydavanych-kvalifikovanych-prostredku-pro-vytvareni-elektronickych-podpisu-v-ceske-republice}
  \item DIA - ověřování zahraničních kvalifikovaných certifikátů: \url{https://www.dia.gov.cz/cs/legislativa/eidas-sluzby-vytvarejici-duveru-a-elektronicka-identifikace/overovani-zahranicnich-kvalifikovanych-certifikatu}
  \item NATO - kolektivní obrana a článek 5: \url{https://www.nato.int/en/what-we-do/introduction-to-nato/collective-defence-and-article-5}
\end{itemize}
